\begin{cuzchapter}{相关技术介绍}{chap:literature}
	% \begin{cuzchapter}{文献综述}{chap:literature}

	%---------------------------------------------------------------------------%
	%->> 文献综述写作指南
	%---------------------------------------------------------------------------%
	% 文献综述是论文的重要组成部分，主要回顾和评述与研究主题相关的已有研究成果
	% 一个完整的文献综述应包含以下几个部分：
	% 1. 研究领域概述：介绍研究领域的基本概念、发展历程和重要性
	% 2.  分析：分析国内外相关研究的主要成果、研究方法和研究趋势
	% 3. 研究评述：评价已有研究的优缺点，指出研究中存在的问题和不足
	% 4. 研究展望：提出未来研究的方向和可能的突破点
	
	\section{React}\label{sec:React}
	React是由FaceBook团队开发的JavaScript库，
	允许开发者使用JSX语法来编写前端页面。
	React会将JSX文件编译为新的JavaScript文件，而JSX可以看作JavaScript的语法拓展\cite{1023559989.nh}。
	可以在JavaScript中编写HTML，这样在编写HTML页面的同时，也可以更方便地处理交互逻辑，
	而非分将这两者分离开来。
	同时其组件化的设计思路可以帮助开发者对组件进行封装以及逻辑复用，
	% 开发时可以将一些大组件拆分为更小的组件，
	如本平台中的前端UI库 
	lucide-react把每一个SVG图标都封装成可定义参数的React组件，
    使开发者可以像使用普通组件那样以声明式的方式放在界面中，
    改变对应参数就可以设置图片的颜色、大小、描边宽度等。
	为了减少性能开销，React设计了虚拟DOM。
	每次页面发生变化的时候，Raect会将虚拟DOM和真实DOM进行对比，
	% 计算并更新的变化的部分，减少渲染真实DOM的资源消耗\cite{1023036236.nh}。
	通过Diff算法计算并更新的变化的部分，减少渲染真实DOM的资源消耗\cite{1023036236.nh}。
	
	
	\section{NestJS}\label{sec:NestJS}
	NodeJS是一个跨平台、开源的服务端JavaScript运行时环境，
	底层使用C++编写。
	它用异步回调的方式处理网络请求和文件操作，
	不会因网络或磁盘阻塞进程，
	极大增加了服务端的资源利用率\cite{1021830702.nh}。
	% 不仅支持HTTP、DNS等网络模块，
	% 而且在许多模块都提供了事件驱动的接口，
	% 以Socket模块为例，它提供了open、message、close等可以注册的事件。
	% 可以帮助开发者更加方便开发服务端。
	
	而NestJS正是基于NodeJS的服务端框架，支持TypeScript和JavaScript，
	融入了面向对象编程、函数式编程等思想\cite{pham2020developing}，
	同时提供了管道、守卫、拦截器等常用功能，
	其中管道可以用于请求参数格式校验和转化，如将字符串的某一部分解析为整数，
	守卫可以用于权限或身份的校验。
	NestJs项目结构中的每个模块有清晰的结构，主要包括controller、service、module等。
	
	\section{MySQL}\label{sec:MySQL}


	MySQL是一个开源的关系型数据库，
	不仅具有性能高、安全性强、运行时占用内存小的优点\cite{XXDL202314058}，
	而且支持跨平台，可以运行在大部分常用的操作系统上，
	C++、Java、Python、TypeScript等常见的
	编程语言都有对应的框架，能够很方便地操作MySQL数据库，
	如本平台就采用typeorm这样的支持TypeScript的ORM（Object Relational Mapping，对象关系映射）框架
	来管理MySQL数据库。



	% WebSocket协议相较于HTTP 协议来说，具有自己的优势，HTTP协议有一个缺陷：
	% 通信只能由客户端发起，也就是说服务器端的变动不能主动由服务器端向客户端推送，
	% 只有在客户端发起请求时，服务器端才能响应，及时性就大打折扣，当然也可以使用长连接，但这是非常浪费带宽资源的

	
	% \section{权限管理}\label{sec:permission-management}

	% \subsection{RBAC模型}

	% 基于角色的访问控制（RBAC）模型是权限管理系统的经典理论。
	% % RBAC模型通过将权限分配给角色，再将角色分配给用户，实现了权限管理的灵活性和可维护性。
	% RBAC模型的思想就是将权限分配给角色，再将角色分配给用户\cite{XTYY201104011rbac,zhang2013rbac}。
	% 这样，权限不再直接赋予用户，而由角色充当两者的媒介，让权限控制更加灵活和简便。

	% 但是，访问控制系统里的用户、角色、权限都是抽象出来的实体概念\cite{wu2002rbacweb}，
	% 当系统规模变大或是系统需求变得复杂，实体的数量增加，
	% 各种关系就变得极其复杂甚至错乱，
	% 所以，在设计平台权限模型时，需要根据具体系统需求，选择并设计合理的RBAC模型。
	% % 在各类信息系统中得到广泛应用。

	% % 在文档协作场景中，RBAC模型需要进行适应性调整以满足细粒度的权限控制需求。
	% % 本平台设计了四级用户角色权限体系（Owner/Editor/Commenter/Viewer），实现了针对幻灯片文档的细粒度访问控制。

	% \subsection{JWT认证技术}

	% JSON Web Token（JWT）是一种用于安全传输信息和传递声明信息的开放标准\cite{fan2016jwt}。
	% JWT由头部、载荷和签名组成，支持跨域认证和单点登录。

	% 对于JWT认证，
	% 客户端得先发起用户登录请求，服务器那边验证用户信息，
	% 验证没问题后，服务器就给客户端签发一个token，客户端拿到后存下来，
	% 等后面想访问服务器资源的时候，需要在请求头里带上这个token\cite{zhou2019jwttoken}。
	% 服务器收到token之后会检查它是不是合法、有没有过期，要是在有效时间内而且用户身份也有效，
	% 就把请求的资源返回给客户端。

	% % \subsection{文档协作中的权限控制}

	% % 文档协作系统中的权限控制需要平衡安全性和可用性。研究表明，过于严格的权限控制会影响协作效率，而过于宽松的控制则可能导致数据安全问题 。本平台通过结合文档可见性控制（公开/私有）和角色权限系统，实现了安全与效率的平衡。



	

	% % \subsection{AI与文档编辑工具的集成}

	% % 将AI能力集成到文档编辑工具中，可以实现智能化创作体验。Grammarly等工具通过AI技术提供语法检查和写作建议。在幻灯片制作领域，AI可以辅助用户设计布局、选择配色方案和生成可视化图表。

	% % \subsection{智能排版与设计建议}

	% % 排版设计是幻灯片制作的重要环节。研究表明，合理的排版可以显著提升信息传达效果。AI技术可以分析内容结构，提供排版优化建议，如字体大小调整、布局平衡和视觉层次设计。

	% \section{版本管理}\label{sec:version-control}

	% \subsection{Git版本管理原理}

	% Git是版本管理系统的典型代表，通过快照方式存储文件状态，支持分支管理、合并操作以及版本追溯\cite{ghodke2024overviewgit}。
	% 然而，Git的工作原理和流程对非技术用户来说学习门槛较高，不仅需要花费较多时间上手，还容易在操作或理解过程中出错。

	% 通常，开发者编写代码的地方称为工作区间，也就是实际工程文件所在的目录，
	% 修改完代码后，需要先将改动添加到一个叫作暂存区的地方，
	% 再往下是本地仓库，它负责将本地代码与远程代码连接起来，
	% 而远程仓库顾名思义就是存放代码的服务器，例如GitHub和Gitee\cite{SXDS202505026}。

	% 整个流程大致如下：在工作区间中修改文件，然后将改动添加到暂存区，再从暂存区提交到本地仓库，最后把本地仓库的内容推送到远程仓库\cite{XMGJ202006028}；
	% 反过来，从远程仓库拉取代码也是沿着这条路径反向进行。

	% \subsection{简化版本管理界面设计}

	% 为降低版本管理工具的使用门槛，许多平台提出了多种直观的界面设计方案，
	% 其中，语雀的版本历史功能为每个文档提供了直观的时间轴界面。
	% 作者打算借鉴这些设计理念，实现简单的版本管理系统，提供版本历史查看、差异对比和一键回滚功能。

	\section{Slidev}\label{sec:slideshow-tools}

	% \subsection{Slidev技术架构与特点}
	% Markdown作为一种轻量级标记语言，因其简洁性和易读性在文档编写中广泛应用。
	% 基于Markdown的幻灯片工具如Reveal.js、Marp、Slidev等，
	% 允许用户使用纯文本编写幻灯片，并通过CSS样式控制视觉效果。

	% 其中，Slidev是专门为开发者设计的现代幻灯片制作工具，
	Slidev是由Anthony Fu开发的现代幻灯片制作工具，
	它基于Vue，把幻灯片制作与Web开发技术结合起来，支持代码高亮、实时预览和自定义主题，
	编写Markdown文档就可以制作幻灯片，
	当然它不仅仅只是支持简单的Markdown语法，也支持Vue、UnoCSS、Shiki、v-motion等工具和技术，
	用户还可以利用它丰富的插件还可以实现许多有趣生动的动画效果，
	如Shiki可以很清晰展示一系列代码步骤的变化过程，
	用户可以使用v-motion很方便地制作幻灯片动画。
	Slidev先将Markdown文档中的“---”作为分页符，同时在第一页中可以对主题、插件进行配置，
	也可以在各个页面对布局等进行配置。
	而后根据该文档中的内容渲染为Web界面，实现幻灯片的效果。
	它支持热加载以及静态资源打包的方式，
	其中热加载方式可以实现根据Markdown文档的内容实时更新渲染幻灯片页面，
	使用命令如“slidev slides.md --port 4090”,
	而静态资源打包则是将Markdown文档构建为一个静态的单页应用，包括JavaScript、CSS、HTML等资源，
	可以部署到服务器中，
	其使用命令如“slidev build slides.md”。
	
	\section{实时协作技术}\label{sec:collaboration-tech}

	% \subsection{CRDT算法与操作转换技术}

	实时协作编辑技术有CRDT（冲突复制数据类型）和OT（操作转换）。
	% 实时协作编辑技术有CRDT（冲突复制数据类型）和OT（操作转换）\cite{SJCM79F193E0431876530855F41675323DC6}。
	CRDT是一种抽象的数据类型，具有明确的接口，
	两个副本接受到相同的更新操作时都会得到相同的结果，
	而且在这种数据类型定义下的操作是可以交换的，且在任何执行顺序下的结果也是相同的\cite{preguicca2018conflict}，
	% 保证多个用户协作时操作具有交换性和幂等性，
	% 达到无冲突的最终一致性\cite{preguicca2018conflict}。
	% CRDT算法通过设计特定的数据结构，保证多用户协作时操作的交换性、结合性和幂等性，从而实现无冲突的最终一致性。
	OT技术则通过转换并发操作来解决冲突，并且需要中心服务器进行协调\cite{takeichi2024coordination}。
	% OT技术则通过转换并发操作\cite{sun1998operational}来解决冲突，并且需要中心服务器进行协调\cite{takeichi2024coordination}。
	% 在现在的Web实时协作领域中，人们常常使用CRDT算法，主要因为它不需要中央节点进行统一调度，且容错能力更强。

	Yjs就是一个使用CRDT算法的高性能实时协作库，内置了文档状态同步和冲突解决机制，
	它使用基于双向链表的CRDT数据结构，
	当某个元素被删除时，并非不会马上从链表中移除，而是将其内容设置为空，
	只有其垃圾回收机制判定不会有冲突时候，才会从链表中移除，
	在保持内存占用较低的同时又获得较好的性能表现\cite{nicolaescu2015yjs}。
	现在，Yjs已经可以应用于CodeMirror和Quill等Web编辑器。

	% \subsection{WebSocket实时通信技术}

	WebSocket协议提供了全双工通信通道，这意味着服务器与浏览器之间可以相互发送和接收消息，
	能够实现浏览器和服务器间的异步通信\cite{1024463893.nh,zhou2019websocket}，这是实现实时协作的关键技术。
	% 但对于HTTP协议，通信只能由客户端发起，服务器的消息却不能主动推送到浏览器，
	% 这样的性能和协作及时性将大打折扣。
	Hocuspocus库使用Yjs，可以基于WebSocket协议实现多客户端的文档状态同步，
	而且它提供了文档初始化、存储、认证等事件钩子，开发者可以很方便地利用这些钩子
	完成自己项目中的业务需求。

	% Slidev（slide + dev）是专门为开发者设计的现代幻灯片制作工具。
	% 它基于Vue.js组件系统，支持Markdown语法编写幻灯片内容，并提供了丰富的插件生态。
	% Slidev的核心理念是将幻灯片制作与Web开发技术结合，支持代码高亮、实时预览和自定义主题。

	% \subsection{基于Markdown的幻灯片制作}


	% \subsection{开发者工具集成}

	% 现代Web幻灯片工具注重与开发者工具的集成。
	% VS Code等代码编辑器提供了Slidev插件支持。
	% 本平台通过集成CodeMirror编辑器，为开发者提供了熟悉的编码环境和实时预览功能。

	\section{Qwen-Plus}
	Qwen-Plus是阿里巴巴团队开发的多模态大语言模型，
	使用成本较低，支持中文、英文等主流语言。
	不仅能理解较长的上下文并可以进行一定的逻辑推理，快速生成连贯的文本内容，
	常常用于生成高质量的文档内容、提供个性化的建议、概括分析文本内容等，
	而且能理解并分析图像、视频。
	开发者可以很方便将阿里云提供的Qwen-Plus接口服务接入到自己的应用中。
	% 不仅，可以处理代码生成、逻辑推理等复杂任务，
	% 而且使用成本较低，能理解较长的上下文并快速生成连贯的文本内容以及进行逻辑推理，
	% 常常用于生成高质量的文档内容、提供个性化的建议、概括分析文本内容等。
	% 常应用于智能客服、舆情分析、个人智能助手等领域。
	% 在本平台中，可以将Slidev的相关文档加载到它的上下文中，
	% 以此来辅助用户快速生成幻灯片文档大纲以及优化部分内容的表达。

	% \section{AI辅助创作}\label{sec:ai-assistance}

	% % \subsection{大语言模型与Slidev的融合}

	% 大语言模型（LLMs）如GPT系列、通义千问等，擅长处理各种自然语言任务。
	% 这些模型通过预训练和微调，能够理解上下文并生成连贯的文本内容。
	% 在文档创作领域，LLMs可以辅助用户生成初稿、优化表达和提供创意建议。

	% 国内也有研究者将大语言模型应用于Slidev的内容生成\cite{slidevAi}，并开发了名为PreGenie的Agent框架，他们将多模态的文档输入到该框架，
	% 生成出Slidev的源码，进而渲染出精美的幻灯片。


	% \section{研究评述与展望}\label{sec:review-prospect}

	% \subsection{现有研究的不足}
	% 通过对实时协作技术、权限管理、版本管理和幻灯片工具等领域的研究，
	% 可以发现现有研究存在一些不足之处。
	% % 实时协作技术的应用多集中于通用文档处理，与幻灯片制作场景的结合还不够深入；
	% 专业的版本管理工具虽然成熟，但学习门槛较高，缺少面向普通用户的简化操作界面；
	% 而Slidev作为一款优秀的幻灯片渲染工具，也缺乏对团队协作的原生支持，难以直接满足多人协同制作的需求。
	% % 通过对实时协作技术、权限管理、AI辅助创作、版本管理和幻灯片工具等领域的文献综述，可以发现以下不足：

	% % 实时协作与幻灯片制作的结合不足：现有研究主要关注通用文档的实时协作，针对幻灯片制作场景的优化研究相对较少。
	% % 权限管理的用户体验研究不足：大多数权限管理系统注重安全性，但对非技术用户的易用性考虑不够充分。
	% % 版本管理工具的学习门槛高：Git等专业版本管理工具对非技术用户存在较高的学习成本，缺乏简化操作界面。
	% % Slidev的协作功能缺失：Slidev作为优秀的单机幻灯片工具，缺乏对团队协作的原生支持。

	% \subsection{未来研究方向}
	% 考虑到上述不足，未来的研究可以关注面向幻灯片场景的协作，
	% 设计直观的权限管理模型，用图形化方式减少权限配置的复杂度，让用户使用更直观和简便；
	% % 探索AI与幻灯片语法的深度集成，开发专门的提示词模板和上下文管理机制，提高AI生成内容与Slidev语法的直接适配性；
	% 简化版本管理操作，用时间轴显示幻灯片文档的历史版本，提供一键回滚、差异可视化等功能；
	% 以及基于Slidev设计开发一个多人幻灯片协作平台，从单机幻灯片渲染工具转变为团队协作平台。
	% 基于上述不足，未来的研究可以从以下几个方向展开：

	% 面向幻灯片场景的协作优化：研究CRDT算法在幻灯片编辑中的特殊优化，如幻灯片页面的插入、删除和重排序操作。
	% 直观的权限管理界面设计：设计图形化权限管理界面，降低权限配置的复杂度，提高用户体验。
	% AI与幻灯片语法的深度集成：开发专门的提示词模板和上下文管理机制，提高AI生成内容与Slidev语法的直接适配性。
	% 简化版本管理操作：设计直观的版本管理界面，提供一键回滚、差异可视化等简化功能。
	% Slidev协作功能扩展：为Slidev添加实时协作支持，使其从单机工具转变为团队协作平台。

	% \subsection{本研究的创新点}
	% 本平台将实时协作、权限管理、AI辅助与版本管理整合进统一的幻灯片制作环境，形成一套AI驱动的协作工作流。
	% 编辑层面采用Yjs的CRDT算法，在支撑多用户无冲突同步编辑的同时，还能追踪实时光标位置。
	% 权限体系基于RBAC模型，划分出Owner、Editor、Commenter和Viewer四种角色，以实现细粒度的访问控制。
	% AI能力通过提示词工程与上下文管理，深度适配Slidev语法，让生成内容可以直接用于幻灯片而无需繁琐调整。
	% 版本管理则着力简化交互界面，使非技术背景的用户也能轻松理解与操作，从而在日常团队协作中整体拉高效率。
	
	% 本平台在以下方面具有创新性：

	% AI驱动的幻灯片协作平台：将实时协作、权限管理、AI辅助和版本管理整合到统一的幻灯片制作平台中。
	% 基于CRDT的实时协作实现：采用Yjs CRDT算法实现无冲突的幻灯片实时编辑，支持多用户同时编辑和实时光标追踪。
	% 四级权限管理系统：基于RBAC模型设计Owner/Editor/Commenter/Viewer四级权限体系，实现细粒度的文档访问控制。
	% AI与Slidev语法的深度集成：通过提示词工程和上下文管理，提高AI生成内容与Slidev语法的直接适配性。
	% 简化版本管理界面：设计直观的版本管理界面，降低非技术用户的使用门槛，提高团队协作效率。

\end{cuzchapter}